% GENERATED FILE. Edit canonical lessons, then run npm run generate:books.
% canonical-source-hash: fnv1a64:93a1b20b9562b178
% canonical-lessons: BN-C121-pyacha, BN-C121-tiya, BN-C121-kobutor, BN-C121-chorui, BN-C121-bok

\chapter{Owl, Parrot, Pigeon, Sparrow, Heron}
\label{ch:bn-owl-parrot-pigeon-sparrow-heron}
\hlchaptermodality{121}

\begin{chapteropening}
\textbf{By the end of this chapter:} \emph{I can say an owl, a parrot, a pigeon, a sparrow and a heron.}

Complete the last lesson of chapter 121: I can say an owl, a parrot, a pigeon, a sparrow and a heron.
\end{chapteropening}

\section[\texorpdfstring{\bn{প্যাঁচা} (pyā̃chā)}{প্যাঁচা (pyā̃chā)}]{\bn{প্যাঁচা} (pyā̃chā) — an owl}

\label{lesson:BN-C121-pyacha}



\begin{quote}
Before the new one: say the Bengali for a jackal, then the Bengali for a bat.
\end{quote}

\subsection*{You'll want to know: \bn{প্যাঁচা}}
\textbf{\bn{প্যাঁচা}} — \emph{pyā̃chā} — \textquotedblleft{}an owl\textquotedblright{}.

\begin{grammarlens}[title={What you've built}]
One more word for this chapter.
\end{grammarlens}

\subsection*{Your turn}
\begin{itemize}
\raggedright
  \item \emph{Say it:} \emph{pyā̃chā}
  \item \emph{Say it:} \emph{pyā̃chā}, once more
  \item \emph{Say it:} say \emph{pyā̃chā}
  \item \emph{Recall:} say the Bengali for a jackal, then the Bengali for a bat, then say \emph{pyā̃chā} again
\end{itemize}

\begin{tcolorbox}[breakable,colback=teal!4,colframe=teal!35!black,title={Before you move on}]
What does \bn{প্যাঁচা} mean? (\textquotedblleft{}An owl\textquotedblright{}.) Say it once more.
\end{tcolorbox}
\section[\texorpdfstring{\bn{টিয়া} (tiyā)}{টিয়া (tiyā)}]{\bn{টিয়া} (tiyā) — a parrot}

\label{lesson:BN-C121-tiya}



\begin{quote}
Before the new one: say the Bengali for a bat, then the Bengali for an owl.
\end{quote}

\subsection*{You'll want to know: \bn{টিয়া}}
\textbf{\bn{টিয়া}} — \emph{tiyā} — \textquotedblleft{}a parrot\textquotedblright{}.

\begin{grammarlens}[title={What you've built}]
One more word for this chapter.
\end{grammarlens}

\subsection*{Your turn}
\begin{itemize}
\raggedright
  \item \emph{Say it:} \emph{tiyā}
  \item \emph{Say it:} \emph{tiyā}, once more
  \item \emph{Say it:} say \emph{tiyā}
  \item \emph{Recall:} say the Bengali for a bat, then the Bengali for an owl, then say \emph{tiyā} again
\end{itemize}

\begin{tcolorbox}[breakable,colback=teal!4,colframe=teal!35!black,title={Before you move on}]
What does \bn{টিয়া} mean? (\textquotedblleft{}A parrot\textquotedblright{}.) Say it once more.
\end{tcolorbox}
\section[\texorpdfstring{\bn{কবুতর} (kôbutôr)}{কবুতর (kôbutôr)}]{\bn{কবুতর} (kôbutôr) — a pigeon}

\label{lesson:BN-C121-kobutor}



\begin{quote}
Before the new one: say the Bengali for an owl, then the Bengali for a parrot.

Say \textquotedblleft{}I live in Kolkata,\textquotedblright{} then \textquotedblleft{}she speaks.\textquotedblright{} (\emph{Āmi Kolkātāy thāki}; \emph{se bôle}, the same for he and she.)
\end{quote}

\subsection*{You'll want to know: \bn{কবুতর}}
\textbf{\bn{কবুতর}} — \emph{kôbutôr} — \textquotedblleft{}a pigeon\textquotedblright{}.

\begin{grammarlens}[title={What you've built}]
One more word for this chapter.
\end{grammarlens}

\subsection*{Your turn}
\begin{itemize}
\raggedright
  \item \emph{Say it:} \emph{kôbutôr}
  \item \emph{Say it:} \emph{kôbutôr}, once more
  \item \emph{Say it:} say \emph{kôbutôr}
  \item \emph{Recall:} say the Bengali for an owl, then the Bengali for a parrot, then say \emph{kôbutôr} again
\end{itemize}

\begin{tcolorbox}[breakable,colback=teal!4,colframe=teal!35!black,title={Before you move on}]
What does \bn{কবুতর} mean? (\textquotedblleft{}A pigeon\textquotedblright{}.) Say it once more.
\end{tcolorbox}
\section[\texorpdfstring{\bn{চড়ুই} (chôṛui)}{চড়ুই (chôṛui)}]{\bn{চড়ুই} (chôṛui) — a sparrow}

\label{lesson:BN-C121-chorui}



\begin{quote}
Before the new one: say the Bengali for a parrot, then the Bengali for a pigeon.
\end{quote}

\subsection*{You'll want to know: \bn{চড়ুই}}
\textbf{\bn{চড়ুই}} — \emph{chôṛui} — \textquotedblleft{}a sparrow\textquotedblright{}.

\begin{grammarlens}[title={What you've built}]
One more word for this chapter.
\end{grammarlens}

\subsection*{Your turn}
\begin{itemize}
\raggedright
  \item \emph{Say it:} \emph{chôṛui}
  \item \emph{Say it:} \emph{chôṛui}, once more
  \item \emph{Say it:} say \emph{chôṛui}
  \item \emph{Recall:} say the Bengali for a parrot, then the Bengali for a pigeon, then say \emph{chôṛui} again
\end{itemize}

\begin{tcolorbox}[breakable,colback=teal!4,colframe=teal!35!black,title={Before you move on}]
What does \bn{চড়ুই} mean? (\textquotedblleft{}A sparrow\textquotedblright{}.) Say it once more.
\end{tcolorbox}
\section[\texorpdfstring{\bn{বক} (bôk)}{বক (bôk)}]{\bn{বক} (bôk) — a heron}

\label{lesson:BN-C121-bok}



\begin{quote}
Before the new one: say the Bengali for a pigeon, then the Bengali for a sparrow.
\end{quote}

\subsection*{You'll want to know: \bn{বক}}
\textbf{\bn{বক}} — \emph{bôk} — \textquotedblleft{}a heron\textquotedblright{}.

\begin{grammarlens}[title={What you've built}]
One more word for this chapter.
\end{grammarlens}

\subsection*{Your turn}
\begin{itemize}
\raggedright
  \item \emph{Say it:} \emph{bôk}
  \item \emph{Say it:} \emph{bôk}, once more
  \item \emph{Say it:} say \emph{bôk}
  \item \emph{Recall:} say the Bengali for a pigeon, then the Bengali for a sparrow, then say \emph{bôk} again
\end{itemize}

\begin{tcolorbox}[breakable,colback=teal!4,colframe=teal!35!black,title={Before you move on}]
What does \bn{বক} mean? (\textquotedblleft{}A heron\textquotedblright{}.) Say it once more.
\end{tcolorbox}
